%!TEX program = xelatex
\PassOptionsToPackage{full}{textcomp}
\documentclass[lang=en,11pt,a4paper,citestyle=authoryear,nobib=true]{elegantpaper}

\title{HW8 - STAT 709}
\author{Wells Guan}

\usepackage{array,url,stix}
\usepackage{tikz}
\usepackage{tikz-cd}
\usepackage{enumitem}
\usepackage{amsmath}
\usepackage{amssymb}
\usepackage{listings}
\usepackage{subfigure}

\lstset{
    frame=single,
    columns=fullflexible,
    keepspaces=true,
    showstringspaces=false,
    breaklines=true,
    breakatwhitespace=false,
    basicstyle=\ttfamily\footnotesize,
    numbers=left,
    numberstyle=\tiny,
    numbersep=6pt
}

\newcommand{\Z}{\mathbb{Z}}
\newcommand{\R}{\mathbb{R}}
\newcommand{\N}{\mathbb{N}}
\newcommand{\B}{\mathcal{B}}
\newcommand{\E}{\mathbb{E}}
\renewcommand{\P}{\mathbb{P}}
\newcommand{\norm}[1]{\left\lVert#1\right\rVert}

\begin{document}

\maketitle

\noindent All three problems are required. Complete every part of each problem.

\subsection*{Problem 1. Fisher information and sufficient statistics}
Jun Shao, \textit{Mathematical Statistics}, second edition, Chapter 3:
Exercise 36, parts (a) and (b) (p.~222).

Let $X$ be a sample having a p.d.f.\ satisfying the conditions in
Theorem 3.3, where $\theta$ is a $k$-vector of unknown parameters, and let
$T(X)$ be a statistic. If $T$ has a p.d.f.\ $g_\theta$ satisfying the
conditions in Theorem 3.3, then we define
\[
I_T(\theta)=E\left\{
\frac{\partial}{\partial\theta}\log g_\theta(T)
\left[\frac{\partial}{\partial\theta}\log g_\theta(T)\right]^\tau
\right\}
\]
to be the Fisher information about $\theta$ contained in $T$.
\begin{enumerate}[label=(\alph*),leftmargin=*]
\item Show that $I_X(\theta)-I_T(\theta)$ is nonnegative definite,
where $I_X(\theta)$ is the Fisher information about $\theta$ contained in $X$.
\item Show that $I_X(\theta)=I_T(\theta)$ if $T$ is sufficient for $\theta$.
\end{enumerate}

\par\vspace{0.5em}\noindent\textbf{Sol.}\par
% Write your solution to Problem 1 here.
\vspace{3em}

\subsection*{Problem 2. Bias--variance tradeoff and the Cram\'er--Rao bound}
Let the sample $X$ have density $f_\theta$ with respect to a fixed
measure, where the scalar parameter $\theta$ ranges over an open interval.
Let $T=T(X)$ be a statistic estimating $\theta$, and write
\[
b_T(\theta)=\E_\theta T-\theta,\qquad
I(\theta)=\E_\theta\left[
\left\{\frac{\partial}{\partial\theta}\log f_\theta(X)\right\}^2
\right].
\]
Assume the usual Cram\'er--Rao regularity conditions: $f_\theta$ is
differentiable in $\theta$, its support does not depend on $\theta$,
and differentiation may pass under the integrals of $f_\theta$ and
$Tf_\theta$. Also assume $0<I(\theta)<\infty$ and
$\E_\theta T^2<\infty$. Prove that
\[
\E_\theta[(T-\theta)^2]\ge
\frac{[1+b'_T(\theta)]^2}{I(\theta)}+b_T(\theta)^2.
\]
\noindent\textit{Remark:} A suitable amount of bias can reduce variance
and lead to a smaller mean squared error. The next exercise gives an example.

\par\vspace{0.5em}\noindent\textbf{Sol.}\par
% Write your solution to Problem 2 here.
\vspace{3em}

\subsection*{Problem 3. Variance estimation}
Let $X_1,\ldots,X_n\overset{\mathrm{i.i.d.}}{\sim}N(\mu,\sigma^2)$,
where $n\ge 2$, $\mu\in\R$, and $\sigma^2>0$ are unknown.
Write $\overline X=n^{-1}\sum_{i=1}^n X_i$ and define
\[
S_n^2=\frac{1}{n-1}\sum_{i=1}^n(X_i-\overline X)^2,\qquad
V_n=\frac1n\sum_{i=1}^n(X_i-\overline X)^2=\frac{n-1}{n}S_n^2.
\]
\begin{enumerate}[label=(\roman*),leftmargin=*]
\item Derive the bias and variance of $V_n$ as an estimator of $\sigma^2$.
Show that its mean squared error is smaller than that of $S_n^2$.
\item Among estimators $c_nS_n^2$, where the nonrandom multiplier $c_n$
may depend on $n$ but not on the data or unknown parameters, find the value
that minimizes mean squared error.
\item Derive the Fisher information matrix for $\theta=(\mu,\sigma^2)^\top$
based on all $n$ observations. The sample mean $\overline X$ and sample
variance $S_n^2$ are the uniformly minimum-variance unbiased estimators
(UMVUEs) of $\mu$ and $\sigma^2$, respectively. Compare the covariance matrix
of $\widehat\theta=(\overline X,S_n^2)^\top$ with the matrix
Cram\'er--Rao bound. State whether each component attains its bound.
\end{enumerate}

\par\vspace{0.5em}\noindent\textbf{Sol.}\par
% Write your solution to Problem 3 here.
\vspace{3em}

\end{document}
